% --- Archon physics formalization source begin ---
% archon:physics
% archon:covers IPhO2026Problems/problem_IPhO_2026_2_C_2.lean
% archon:source-report reports/ipho_2026/problem_IPhO_2026_2_C_2.source.json
% archon:problem-id IPhO_2026_2
% archon:part-id C.2
% archon:previous-part IPhO_2026_2_C_1
% archon:previous-part-policy natural_language_prerequisite_only

\chapter{Physics problem IPhO\_2026\_2\_C\_2}
\label{ch:IPhO2026Problems_problem_IPhO_2026_2_C_2}

\paragraph{Problem source.}
For the half-cylindrical mirror of radius R, ray A is incident at angle theta
and its reflected line is y = m\_A*x + b\_A.  A neighboring parallel ray B is
incident at theta + Delta theta, with Delta theta much smaller than theta, and
its reflected line is y = m\_B*x + b\_B.  The envelope/intersection of neighboring
rays forms the caustic.  Use Figure 2g and its coordinate convention.

Current subquestion:
Expand m\_B and b\_B to first order in Delta theta.

\paragraph{Current subquestion.}
Expand m\_B and b\_B to first order in Delta theta.

\paragraph{Recorded answer/context.}
m\_B = cot(2*theta) - 2*csc(2*theta)\textasciicircum{}2*Delta theta; b\_B = [R/(2*cos(theta))]*(1 + tan(theta)*Delta theta), up to O(Delta theta\textasciicircum{}2).

\paragraph{Figure/image path.}
/root/proposal\_for\_physic/science-mango/ipho\_2026\_source/image/T2\_page-4.png

\paragraph{Reusable previous-part conclusions.}
\begin{itemize}
\item Source C.1. Question: Write the slope m\_A and intercept b\_A of reflected ray A in terms of theta and R. Reusable conclusions: m\_A = cot(2*theta), and b\_A = R/(2*cos(theta)). Policy: natural\_language\_prerequisite\_only; do\_not\_import\_Lean\_output
\end{itemize}

\paragraph{Formalization target.}
create a compiling Lean file with sorry bodies at `IPhO2026Problems/problem\_IPhO\_2026\_2\_C\_2.lean`. The Lean declarations must preserve the physical quantities, dimensions or dimensional roles, figure labels, governing-law hypotheses, and final relation expressed by this problem.
Use Mathlib/Physlib names found through LeanExplore where available. If a physics API is missing, introduce faithful local abstractions rather than scalar placeholder aliases.

\paragraph{Iteration 002 redraft contract.}
Use the same Physlib-backed length carrier and named SI projection as C.1 for
the mirror radius, coordinates, and intercept.  The slope and angular
increment are dimensionless.  Keep the conclusion as two genuine local
first-order expansions with quadratic remainders; for the intercept, apply
the asymptotic statement to its SI-length projection rather than replacing a
length by a bare real.  Retain the nonzero sine/cosine conditions and the
Figure 2g incoming/outgoing orientation.

\begin{theorem}[Physics formalization target]
  \label{thm:physics:IPhO_2026_2_C_2:target}
  \lean{IPhO2026Problems.IPhO2026_2_C_2.rayB_firstOrderExpansion}
  \uses{def:physics:IPhO_2026_2_C_2:length_quantity, def:physics:IPhO_2026_2_C_2:length_si, def:physics:IPhO_2026_2_C_2:aux001, def:physics:IPhO_2026_2_C_2:aux002, def:physics:IPhO_2026_2_C_2:aux003, def:physics:IPhO_2026_2_C_2:aux004, def:physics:IPhO_2026_2_C_2:aux005, def:physics:IPhO_2026_2_C_2:aux006, def:physics:IPhO_2026_2_C_2:aux007, def:physics:IPhO_2026_2_C_2:aux008, def:physics:IPhO_2026_2_C_2:aux009, def:physics:IPhO_2026_2_C_2:aux010, def:physics:IPhO_2026_2_C_2:aux011, def:physics:IPhO_2026_2_C_2:aux012, def:physics:IPhO_2026_2_C_2:aux013, lem:physics:IPhO_2026_2_C_2:aux014, lem:physics:IPhO_2026_2_C_2:aux015}
  IPhO 2026 Problem 2 C.2: the reflected line of the neighboring ray B has the stated first-order slope and intercept expansions as Δθ → 0. The two IsBigO conclusions say that the displayed remainders are bounded by a constant multiple of Δθ² near zero. Thus the approximation order in the source is part of the theorem contract rather than being silently discarded.
\end{theorem}
\begin{proof}
  Apply the exact C.1 coefficient laws at \(\theta+\Delta\theta\).  The
  derivatives of \(\cot(2\theta)\) and \(R/(2\cos\theta)\) are respectively
  \(-2\csc^2(2\theta)\) and
  \(R\tan\theta/(2\cos\theta)\).  Taylor expansion on a neighborhood where
  the stated denominators stay nonzero leaves remainders bounded by a constant
  times \((\Delta\theta)^2\).
\end{proof}
\section*{Formal declaration index}
The following blocks record the typed carriers and bridge declarations used by the target.  Their dependency lists follow the declaration-level model in the covered Lean file.

\begin{definition}[Declaration LengthQuantity]
  \label{def:physics:IPhO_2026_2_C_2:length_quantity}
  \lean{IPhO2026Problems.IPhO2026_2_C_2.LengthQuantity}
  A unit-independent physical length represented by Physlib's dimensionful
  carrier with length dimension.
\end{definition}
\begin{proof}
  This is the specialization of Physlib's dimensionful-quantity interface to
  the physical dimension of length.
\end{proof}

\begin{definition}[Declaration lengthSI]
  \label{def:physics:IPhO_2026_2_C_2:length_si}
  \lean{IPhO2026Problems.IPhO2026_2_C_2.lengthSI}
  \uses{def:physics:IPhO_2026_2_C_2:length_quantity}
  The scalar projection of a physical length in metres under Physlib's SI
  unit choice.
\end{definition}
\begin{proof}
  Evaluate the dimensionful length at the SI unit choice and take its
  underlying real readout.
\end{proof}

\begin{definition}[Declaration Point2D]
  \label{def:physics:IPhO_2026_2_C_2:aux001}
  \lean{IPhO2026Problems.IPhO2026_2_C_2.Point2D}
  \uses{def:physics:IPhO_2026_2_C_2:length_quantity}
  A point in the two-dimensional cross-section used in Figure 2g.  Both
  coordinates are Physlib length quantities with a shared named SI projection.
\end{definition}
\begin{proof}
  This is the typed definition of the stated physical carrier; its fields or defining equation provide the claimed data.
\end{proof}

\begin{definition}[Declaration Direction2D]
  \label{def:physics:IPhO_2026_2_C_2:aux002}
  \lean{IPhO2026Problems.IPhO2026_2_C_2.Direction2D}
  A directed vector in the Figure 2g coordinate frame. Its components are used only to retain the incoming/outgoing orientation of an optical ray.
\end{definition}
\begin{proof}
  This is the typed definition of the stated physical carrier; its fields or defining equation provide the claimed data.
\end{proof}

\begin{definition}[Declaration verticalIncomingDirection]
  \label{def:physics:IPhO_2026_2_C_2:aux003}
  \lean{IPhO2026Problems.IPhO2026_2_C_2.verticalIncomingDirection}
  \uses{def:physics:IPhO_2026_2_C_2:aux002}
  The upward vertical direction of the parallel incident rays in Figure 2g.
\end{definition}
\begin{proof}
  This is the typed definition of the stated physical carrier; its fields or defining equation provide the claimed data.
\end{proof}

\begin{definition}[Declaration AffineLineReadout]
  \label{def:physics:IPhO_2026_2_C_2:aux004}
  \lean{IPhO2026Problems.IPhO2026_2_C_2.AffineLineReadout}
  \uses{def:physics:IPhO_2026_2_C_2:length_quantity, def:physics:IPhO_2026_2_C_2:length_si}
  A reflected straight line with dimensionless slope and Physlib
  length-valued intercept.  Its scalar analytic equation uses the intercept's
  named SI projection.
\end{definition}
\begin{proof}
  This is the typed definition of the stated physical carrier; its fields or defining equation provide the claimed data.
\end{proof}

\begin{definition}[Declaration Figure2gRayLabel]
  \label{def:physics:IPhO_2026_2_C_2:aux005}
  \lean{IPhO2026Problems.IPhO2026_2_C_2.Figure2gRayLabel}
  Incidence-point labels from Figure 2g.
\end{definition}
\begin{proof}
  This is the typed definition of the stated physical carrier; its fields or defining equation provide the claimed data.
\end{proof}

\begin{definition}[Declaration OpticalRay2D]
  \label{def:physics:IPhO_2026_2_C_2:aux006}
  \lean{IPhO2026Problems.IPhO2026_2_C_2.OpticalRay2D}
  \uses{def:physics:IPhO_2026_2_C_2:aux001, def:physics:IPhO_2026_2_C_2:aux002, def:physics:IPhO_2026_2_C_2:aux004, def:physics:IPhO_2026_2_C_2:aux005}
  The data retained for an incident ray and its reflected branch. Angles are real numbers measured in radians in the convention of Figure 2g.
\end{definition}
\begin{proof}
  This is the typed definition of the stated physical carrier; its fields or defining equation provide the claimed data.
\end{proof}

\begin{definition}[Declaration HalfCylindricalMirror]
  \label{def:physics:IPhO_2026_2_C_2:aux007}
  \lean{IPhO2026Problems.IPhO2026_2_C_2.HalfCylindricalMirror}
  \uses{def:physics:IPhO_2026_2_C_2:length_quantity, def:physics:IPhO_2026_2_C_2:length_si}
  The half-cylindrical mirror, represented by a positive Physlib
  length-valued circular cross-section radius.
\end{definition}
\begin{proof}
  This is the typed definition of the stated physical carrier; its fields or defining equation provide the claimed data.
\end{proof}

\begin{definition}[Declaration OnUpperSemicircle]
  \label{def:physics:IPhO_2026_2_C_2:aux008}
  \lean{IPhO2026Problems.IPhO2026_2_C_2.OnUpperSemicircle}
  \uses{def:physics:IPhO_2026_2_C_2:aux001, def:physics:IPhO_2026_2_C_2:aux007}
  A point lies on the upper semicircular cross-section of the mirror.
\end{definition}
\begin{proof}
  This is the typed definition of the stated physical carrier; its fields or defining equation provide the claimed data.
\end{proof}

\begin{definition}[Declaration AffineLineReadout.Contains]
  \label{def:physics:IPhO_2026_2_C_2:aux009}
  \lean{IPhO2026Problems.IPhO2026_2_C_2.AffineLineReadout.Contains}
  \uses{def:physics:IPhO_2026_2_C_2:aux001, def:physics:IPhO_2026_2_C_2:aux004}
  A point lies on a reflected line in the signed coordinate convention of Figure 2g.
\end{definition}
\begin{proof}
  This is the typed definition of the stated physical carrier; its fields or defining equation provide the claimed data.
\end{proof}

\begin{definition}[Declaration HasFigure2gGeometry]
  \label{def:physics:IPhO_2026_2_C_2:aux010}
  \lean{IPhO2026Problems.IPhO2026_2_C_2.HasFigure2gGeometry}
  \uses{def:physics:IPhO_2026_2_C_2:aux003, def:physics:IPhO_2026_2_C_2:aux005, def:physics:IPhO_2026_2_C_2:aux006, def:physics:IPhO_2026_2_C_2:aux007, def:physics:IPhO_2026_2_C_2:aux008, def:physics:IPhO_2026_2_C_2:aux009}
  The geometry and orientation read directly from Figure 2g for a ray whose incidence angle is φ. In particular, the incoming ray is vertical and the reflected ray is taken on the left-going branch.
\end{definition}
\begin{proof}
  This is the typed definition of the stated physical carrier; its fields or defining equation provide the claimed data.
\end{proof}

\begin{definition}[Declaration HaveParallelIncomingDirections]
  \label{def:physics:IPhO_2026_2_C_2:aux011}
  \lean{IPhO2026Problems.IPhO2026_2_C_2.HaveParallelIncomingDirections}
  \uses{def:physics:IPhO_2026_2_C_2:aux006}
  The two incident rays are parallel with the same directed incoming orientation, as specified in C.2.
\end{definition}
\begin{proof}
  This is the typed definition of the stated physical carrier; its fields or defining equation provide the claimed data.
\end{proof}

\begin{definition}[Declaration SatisfiesPreviousPartC1]
  \label{def:physics:IPhO_2026_2_C_2:aux012}
  \lean{IPhO2026Problems.IPhO2026_2_C_2.SatisfiesPreviousPartC1}
  \uses{def:physics:IPhO_2026_2_C_2:aux006, def:physics:IPhO_2026_2_C_2:aux007}
  The reusable result of C.1 for the ray labelled A. This is previous-part input, rather than either of the first-order conclusions requested in C.2.
\end{definition}
\begin{proof}
  This is the typed definition of the stated physical carrier; its fields or defining equation provide the claimed data.
\end{proof}

\begin{definition}[Declaration SatisfiesHalfCylindricalSpecularLaw]
  \label{def:physics:IPhO_2026_2_C_2:aux013}
  \lean{IPhO2026Problems.IPhO2026_2_C_2.SatisfiesHalfCylindricalSpecularLaw}
  \uses{def:physics:IPhO_2026_2_C_2:aux006, def:physics:IPhO_2026_2_C_2:aux007}
  The exact coefficient consequence of specular reflection from the circular mirror at an arbitrary incidence angle. C.2 applies this governing law at θ + Δθ and then takes its first-order expansion.
\end{definition}
\begin{proof}
  This is the typed definition of the stated physical carrier; its fields or defining equation provide the claimed data.
\end{proof}

\begin{lemma}[Declaration slope\_eq\_of\_specular\_law]
  \label{lem:physics:IPhO_2026_2_C_2:aux014}
  \lean{IPhO2026Problems.IPhO2026_2_C_2.slope_eq_of_specular_law}
  \uses{def:physics:IPhO_2026_2_C_2:aux006, def:physics:IPhO_2026_2_C_2:aux007, def:physics:IPhO_2026_2_C_2:aux013}
  The exact slope equation exposed by the circular-mirror reflection law.
\end{lemma}
\begin{proof}
  Substitute the cited governing relations in order and use the stated positivity, orientation, and branch conditions to obtain the conclusion.
\end{proof}

\begin{lemma}[Declaration intercept\_eq\_of\_specular\_law]
  \label{lem:physics:IPhO_2026_2_C_2:aux015}
  \lean{IPhO2026Problems.IPhO2026_2_C_2.intercept_eq_of_specular_law}
  \uses{def:physics:IPhO_2026_2_C_2:aux006, def:physics:IPhO_2026_2_C_2:aux007, def:physics:IPhO_2026_2_C_2:aux013, lem:physics:IPhO_2026_2_C_2:aux014}
  The exact intercept equation exposed by the circular-mirror reflection law.
\end{lemma}
\begin{proof}
  Substitute the cited governing relations in order and use the stated positivity, orientation, and branch conditions to obtain the conclusion.
\end{proof}
% --- Archon physics formalization source end ---
